% ============================================================
% Section 8: Conclusion
% ============================================================
% Target length: ~0.3 column (150–200 words)
% Write: Month 7, after evaluation is complete
% ============================================================

\section{Conclusion}
\label{sec:conclusion}

\todo{Write after evaluation data is in.
      Structure:
      Para 1: restate the problem and our approach in 2 sentences
      Para 2: summarise the 3 key results (enforcement, latency, bandwidth)
      Para 3: the broader point -- architectural enforcement as a
              necessary complement to model-level defenses
      Para 4: future work (2--3 bullet points max)}

% TEMPLATE:
%
% We presented \ferrite{}, the first browser to enforce capability-based
% governance over AI agent and extension actions with cryptographically
% verifiable audit trails.
% Our evaluation demonstrates that \ferrite{} blocks 100\% of
% \todo{N} adversarial scenarios including reproductions of four
% real-world 2025 attacks, with median broker latency of \todo{X}$\mu$s.
%
% When OpenAI reports that prompt injection ``may never be fully solved''
% at the model level, \ferrite{}'s capability broker provides a
% complementary architectural layer: even a successfully compromised
% agent cannot exceed its granted capabilities, and every decision
% is verifiably recorded.
%
% Future work includes formal verification of capability token properties,
% integration with public transparency logs (Sigstore Rekor), and
% evaluation on the CEF/Chromium Track B to validate engine-agnostic deployment.
